Define Natural Transformations

Given two functors \(F, G: \mathcal{C} \to \mathcal{D}\), a natural transformation \(\eta\) from \(F\) to \(G\), 
denoted as \(\eta: F \Rightarrow G\), is a family of morphisms \(\eta_X: FX \to GX\) for each object \(X \in \mathcal{C}\) 
such that the following diagram commutes for every morphism \(f: X \to Y\) in \(\mathcal{C}\):
\[
\begin{array}{ccc}
FX & \xrightarrow{\eta_X} & GX \\
Ff \downarrow & & \downarrow Gf \\
FY & \xrightarrow{\eta_Y} & GY \\
\end{array}
\]
For any morphism \(f: X \to Y\), 
\(\eta_Y \circ Ff = Gf \circ \eta_X\)
Let \( Z \in \in \mathcal{C} \)
The \( \eta_Z \) are also referred to as components at \( Z \) respectively.
If for every \( Z \), \( \eta_Z \) is an isomorphism in the category \( \mathcal{D} \)
then \( \eta \) is a natural isomorphism. 
More generally any pair \( F : \mathcal{C} \to \mathcal{D} \)
and \( F' : \mathcal{D} \to \mathcal{C} \) s.t. \( F \circ F' = \text{id}_{\mathcal{C}} \) and 
\( F' \circ F = \text{id}_{\mathcal{D}} \) is defined as a natural isomorphism or equivalence of categories.

In the category of functors \( \mathcal{F}(\mathcal{C}, \mathcal{D}) \) natural transformations are precisely the morphisms.


Define fibered products

Suppose we have morphisms \(\alpha: X \to Z\) and \(\beta: Y \to Z\) (in any category). Then the fibered product (or fibred product is an object \(X \times_Z Y\) 
along with morphisms \(\mathrm{pr}_X: X \times_Z Y \to X\) and \(\mathrm{pr}_Y: X \times_Z Y \to Y\), where the two compositions 
\(\alpha \circ \mathrm{pr}_X\), \(\beta \circ \mathrm{pr}_Y: X \times_Z Y \to Z\) agree, such that given any object \(W\) with maps to \(X\) and \(Y\) 
(whose compositions to \(Z\) agree), these maps factor through some unique \(W \to X \times_Z Y\):

INSERT DIAGRAM 39

(Warning: the definition of the fibered product depends on \(\alpha\) and \(\beta\), even though they are omitted from the notation \(X \times_Z Y\).)

By the usual universal property argument, if it exists, it is unique up to unique isomorphism. (You should think this through until it is clear to you.)
As an example, if \(\pi: X \to Y\) is a morphism, and the fibered product \(X \times_Y X\) exists, 
then this determines a diagonal morphism \(\delta_\pi: X \to X \times_Y X\). The diagonal morphism will turn out to be a very useful notion.

Depending on your religion, the diagram

INSERT DIAGRAM 40

is called a fibered/pullback/Cartesian diagram/square (six possibilities — and even more are possible if you prefer “fibred” to “fibered”).
The right way to interpret the notion of fibered product is first to think about what it means in the category of sets.
In Sets it holds
\[X \times_Z Y = \{(x, y) \in X \times Y : \alpha(x) = \beta(y)\}.\]
More precisely, show that the right side, equipped with its evident maps to \(X\) and \(Y\), satisfies the universal property of the fibered product. 
(This will help you build intuition for fibered products.)

The Diagonal Base Change Diagram
Suppose we are given morphisms \(X_1, X_2 \to Y\) and \(Y \to Z\). The following diagram is a Cartesian square.
INSERT IMAGE 41

Two covariant functors \( F : \mathcal{A} \to \mathcal{B} \) and \(G: \mathcal{B} \to \mathcal{A}\) 
are adjoint if there is a natural bijection for all \(A \in \mathcal{A}\) and \(B \in \mathcal{B}\) 
\(\tau_{A,B} : \mathrm{Mor}_{\mathcal{B}}(F(A), B) \to \mathrm{Mor}_{\mathcal{A}}(A, G(B)).\)

We say that \((F, G)\) form an adjoint pair, and that \(F\) is left-adjoint to \(G\) (and \(G\) is right-adjoint to \(F\)). We say \(F\) is a left adjoint 
(and \(G\) is a right adjoint). By “natural” we mean the following. For all \(f: A \to A'\) in \(\mathcal{A}\), we require
INSERT IMAGE 49 

to commute, and for all \(g: B \to B' \) in \(\mathcal{B}\) we want a similar commutative diagram to commute. 
(Here \(f^*\) is the map induced by \(f: A \to A'\), and \(\mathrm{Ff}^*\) is the map induced by \(\mathrm{F}f: \mathrm{F}(A) \to \mathrm{F}(A')\).)
The map \(\tau_{AB} has the following properties. For each \(A\) there is a map \(\eta_A: A \to GF(A)\) so that for any \(g: F(A) \to B\), 
the corresponding \(\tau_{AB}(g): A \to G(B)\) is given by the composition 
\(A \xrightarrow{\eta_A} GF(A) \xrightarrow{Gg} G(B).\)
Similarly, there is a map \(\epsilon_B: FG(B) \to B\) for each \(B\) so that for any \(f: A \to G(B)\), the corresponding map \(\tau_{AB}^{-1}(f): F(A) \to B\) 
is given by the composition 
\(F(A) \xrightarrow{Ff} FG(B) \xrightarrow{\epsilon_B} B.\)
Here is a key example of an adjoint pair.

Suppose \(M, N,\) and \(P\) are \(A\)-modules (where \(A\) is a ring). There's a bijection 
\(\mathrm{Hom}_A(M \otimes_A N, P) \leftrightarrow \mathrm{Hom}_A(M, \mathrm{Hom}_A(N, P))\).

Tensor-Hom Adjunction
\((\cdot) \otimes_A N\) and \(\mathrm{Hom}_A(N, \cdot)\) are adjoint functors.

\subsection*{1.4.B. Exercise}

Show that the map \(\tau_{AB} (1.4.0.1)\) has the following properties. For each \(A\) there is a map \(\eta_A: A \to GF(A)\) so that for any \(g: F(A) \to B\), the corresponding \(\tau_{AB}(g): A \to G(B)\) is given by the composition

\[
A \xrightarrow{\eta_A} GF(A) \xrightarrow{Gg} G(B).
\]

Similarly, there is a map \(\epsilon_B: FG(B) \to B\) for each \(B\) so that for any \(f: A \to G(B)\), the corresponding map \(\tau_{AB}^{-1}(f): F(A) \to B\) is given by the composition

\[
F(A) \xrightarrow{Ff} FG(B) \xrightarrow{\epsilon_B} B.
\]

Here is a key example of an adjoint pair.

Suppose \(M, N,\) and \(P\) are \(A\)-modules (where \(A\) is a ring). 
Describe a bijection \(\mathrm{Hom}_A(M \otimes_A N, P) \leftrightarrow \mathrm{Hom}_A(M, \mathrm{Hom}_A(N, P))\). 
(Hint: try to use the universal property of \(\otimes\).)
Tensor-Hom Adjunction
Show that \((\cdot) \otimes_A N\) and \(\mathrm{Hom}_A(N, \cdot)\) are adjoint functors.


Define additive/abelian categories and sub/quotient objects

Definition
A category \(\mathcal{C}\) is said to be \textit{additive} if it satisfies the following properties.
- For each \(A, B \in \mathcal{C}\), \(\mathrm{Mor}(A, B)\) is an abelian group, such that composition of morphisms distributes over addition. (You should think about what this means — it translates to two distinct statements.)
- \(\mathcal{C}\) has a zero object, denoted \(0\). (This is an object that is simultaneously an initial object and a final object, Definition 1.2.2.)
- It has products of two objects (a product \(A \times B\) for any pair of objects), and hence by induction, products of any finite number of objects.

In an additive category, the morphisms are often called homomorphisms, and \(\mathrm{Mor}\) is denoted by \(\mathrm{Hom}\). In fact, this notation \(\mathrm{Hom}\) is a good indication that you’re working in an additive category. A functor between additive categories preserving the additive structure of \(\mathrm{Hom}\), is called an additive functor.

By the universal property of product and coproduct, we have natural morphisms \(\varphi_{AB}: A \coprod B \to A \times B\). A2. \(\varphi_{AB}\) is an isomorphism. 
This allows us to define a binary operation on \(\mathrm{Mor}(A, B)\), with \(f + g\) (for \(f, g \in \mathrm{Mor}(A, B)\)) defined by the composition 
\(A \xrightarrow{(f,g)} B \times B \xrightarrow{\varphi^{-1}_{B \coprod B}} B \coprod B \to B\)
where the last map is the “codiagonal” defined by universal property of coproduct. 
A little work shows that this endows \(\mathrm{Mor}(A, B)\) with the structure of a \textit{commutative monoid}, i.e., an abelian semigroup with identity. 
The identity is the composition \(A \to 0 \to B\). A3. 
This gives the alternative definition
A0 \( \mathcal{C} \) has a \( 0 \) object
A1 \( \mathcal{C} \) has products and coproducts of any two objects
A2 This commutative monoid \(\mathrm{Mor}(A, B)\) is an abelian group.

One can then define
Given some \( \mathcal{C} \) with a \( 0 \)-object. For a morphism \( f : B \to C \), a kernel is an object \( A \) with morphism \( i : A \to B \)
s.t. \( f \circ i = 0 \). The cokernel reverses these morphisms.
An abelian category is an additive category with
(1) Every map has a kernel and cokernel.
(2) Every monomorphism is the kernel of its cokernel.
(3) Every epimorphism is the cokernel of its kernel.

If \( i : A \to B \) is a monomorphism then we say that \( A \) is a subobject of \( B \),
a quotient object is the dual of this.


